\documentclass[11pt,a4paper]{article}

\usepackage[utf8]{inputenc}
\usepackage[T1]{fontenc}
\usepackage{helvet}
\renewcommand{\familydefault}{\sfdefault}
\usepackage[margin=2.2cm, top=2cm, bottom=2cm]{geometry}
\usepackage{graphicx}
\usepackage{xcolor}
\usepackage{booktabs}
\usepackage{float}
\usepackage{caption}
\usepackage{amsmath}
\usepackage{hyperref}
\usepackage{fancyhdr}
\usepackage{tikz}
\usepackage{enumitem}
\usepackage{tabularx}

\definecolor{bigorange}{RGB}{255,130,0}
\definecolor{biggrey}{RGB}{60,60,60}

\hypersetup{colorlinks=true, linkcolor=bigorange, urlcolor=bigorange}

\pagestyle{fancy}
\fancyhf{}
\renewcommand{\headrulewidth}{0.5pt}
\renewcommand{\headrule}{\hbox to\headwidth{\color{bigorange}\leaders\hrule height \headrulewidth\hfill}}
\fancyhead[L]{\small\color{biggrey} Banco BIG $|$ Rates Research}
\fancyhead[R]{\small\color{biggrey} May 2026}
\fancyfoot[C]{\small\color{biggrey}\thepage}

\begin{document}

% COVER
\thispagestyle{empty}
\begin{tikzpicture}[remember picture, overlay]
    \fill[bigorange] (current page.north west) rectangle ([yshift=-9cm]current page.north east);
    \fill[white, opacity=0.15] ([yshift=-3cm]current page.north west) -- ([yshift=-7cm]current page.north west) -- ([yshift=-5cm]current page.north east) -- ([yshift=-1cm]current page.north east) -- cycle;
    \fill[bigorange] ([yshift=1.5cm]current page.south west) rectangle (current page.south east);
\end{tikzpicture}

\vspace*{1cm}
\begin{center}
    
\includegraphics[width=4.5cm]{../Logo-BIG.png}
\end{center}

\vspace{3.5cm}

\begin{center}
    {\Huge\bfseries\color{biggrey} Interest Rate Spillovers}\\[0.4cm]
    {\Huge\bfseries\color{biggrey} to the Canadian Government Curve}\\[1.0cm]
    {\Large\color{biggrey} 5-Country VAR: US, EU, CA, BR, UK}\\[0.3cm]
    {\Large\color{biggrey} Nelson-Siegel Level, Slope, Curvature}\\[2.0cm]
    {\large\color{biggrey} Rates Research $|$ 5 May 2026}
\end{center}

\vfill
\begin{tikzpicture}[remember picture, overlay]
    \node[anchor=south, yshift=2.0cm] at (current page.south) {
        \color{white}\small\bfseries Banco BIG $|$ Institutional Research
    };
\end{tikzpicture}

\newpage

% PAGE 2: METHODOLOGY & RESULTS
\section*{Methodology: 5-Country VAR on Nelson-Siegel Factors}

We decompose five government yield curves (US, EU, CA, BR, UK) into Level, Slope, and Curvature via Nelson-Siegel ($\lambda=0.10$). A 15-variable VAR(1) is estimated on weekly data (2005--2026, $N=1112$ obs).

\smallskip
\noindent\textbf{Cholesky ordering:} US $\rightarrow$ EU $\rightarrow$ BR $\rightarrow$ UK $\rightarrow$ CA

\smallskip
\noindent Rationale: CA is placed last as the ``receiver'' of interest. US is most exogenous (global reserve currency).

\subsection*{Granger Causality Tests ($\alpha=5\%$)}

\begin{table}[H]
\centering\small
\begin{tabular}{llccc}
\toprule
\textbf{Source} & \textbf{CA Factor} & \textbf{F-stat} & \textbf{p-value} & \textbf{Sig.?} \\
\midrule
    US & Level & 2.95 & 0.032 & \textbf{Sig.} \\
    US & Slope & 1.75 & 0.154 & Not sig. \\
    US & Curvature & 1.23 & 0.297 & Not sig. \\
    EU & Level & 2.43 & 0.063 & Not sig. \\
    EU & Slope & 4.65 & 0.003 & \textbf{Sig.} \\
    EU & Curvature & 1.32 & 0.267 & Not sig. \\
    BR & Level & 9.78 & 0.000 & \textbf{Sig.} \\
    BR & Slope & 3.88 & 0.009 & \textbf{Sig.} \\
    BR & Curvature & 3.32 & 0.019 & \textbf{Sig.} \\
    UK & Level & 0.87 & 0.458 & Not sig. \\
    UK & Slope & 1.96 & 0.118 & Not sig. \\
    UK & Curvature & 0.94 & 0.420 & Not sig. \\
\bottomrule
\end{tabular}
\end{table}

\begin{figure}[H]
    \centering
    \includegraphics[width=\textwidth]{fig_factors.png}
    \vspace{-0.3cm}
    \caption{\small Nelson-Siegel factors for all 5 countries (2005--2026). CA highlighted as the target curve.}
\end{figure}

\newpage

% PAGE 3: IRFs AND FEVD
\section*{Shock Propagation to Canada}

\begin{figure}[H]
    \centering
    \includegraphics[width=\textwidth]{fig_irf.png}
    \vspace{-0.3cm}
    \caption{\small Orthogonalised IRFs: CA factor responses to foreign shocks over 40 weeks. Blue shading = 95\% bootstrap CI.}
\end{figure}

\subsection*{FEVD at 52-Week Horizon (\% of variance)}

\begin{table}[H]
\centering\small
\begin{tabular}{l c c c c c}
\toprule
 & \textbf{US} & \textbf{EU} & \textbf{CA} & \textbf{BR} & \textbf{UK} \\
\midrule
    CA Level & 28.9 & 3.9 & 50.0 & 14.8 & 2.4 \\
    CA Slope & 46.6 & 12.3 & 38.1 & 0.3 & 2.8 \\
    CA Curvature & 49.4 & 4.9 & 38.4 & 6.2 & 1.1 \\
\bottomrule
\end{tabular}
\end{table}

\begin{figure}[H]
    \centering
    \includegraphics[width=\textwidth]{fig_fevd.png}
    \vspace{-0.3cm}
    \caption{\small FEVD: share of CA factor forecast error variance by source country at horizons 1--52 weeks.}
\end{figure}

\newpage

% PAGE 4: MARKET CONTEXT
\section*{Canada: Market Context (May 2026)}


\subsection*{Macro \& Fiscal Backdrop}
\begin{itemize}[nosep, leftmargin=1.2em]
    \item \textbf{Inflation:} CPI at 2.8\% (Mar 2026). Shelter costs remain the largest contributor.
    \item \textbf{Growth:} GDP grew 0.4\% in Q4 2025. Labour market softening; unemployment at 6.7\%.
    \item \textbf{BoC:} Policy rate at 3.25\%, on hold. Markets pricing 1--2 cuts later in 2026.
    \item \textbf{Housing:} Canadian housing vulnerability adds sensitivity to rate movements.
    \item \textbf{Trade:} US tariff uncertainty weighing on business investment.
\end{itemize}

\subsection*{Key Drivers Ahead}
\begin{enumerate}[nosep, leftmargin=1.4em]
    \item \textbf{US dominance:} US explains $\sim$29\% of CA Level and $\sim$47--49\% of Slope and Curvature. The BoC--Fed policy link is the single largest external driver.
    \item \textbf{Brazil connection:} Surprisingly, BR explains $\sim$15\% of CA Level---likely via commodity/EM risk channels.
    \item \textbf{BoC vs Fed divergence:} CAD depreciation constrains BoC's ability to cut independently.
    \item \textbf{Housing correction risk:} Could force BoC to cut aggressively, decoupling from Fed.
\end{enumerate}


\vspace{0.3cm}
\noindent\rule{\textwidth}{0.4pt}
{\footnotesize\color{biggrey} \textbf{Data:} US: FRED (constant-maturity Treasury rates), EU: ECB Statistical Data Warehouse (AAA-rated), CA: Bank of Canada Valet API (benchmark bonds), BR: IPEA/ANBIMA (DI pre-fixed yield curve), UK: Bank of England IADB (nominal spot curve). Model: VAR(1) on weekly NS factors, sample 2005--2026 (1112 obs). This document is for informational purposes only and does not constitute investment advice.}

\end{document}
